% ============================================================
%  OBASI N'DALA - Scheda del personaggio (classe dndbook)
%  Mago - Ordine degli Scriba (livello 5)
%  Umano Variante (PHB 2014 - dati da 5e.tools)
%  Compila con:  pdflatex obasindala.tex   (2 passate)
% ============================================================
\documentclass[letterpaper,twocolumn,openany,nodeprecatedcode]{dndbook}
\usepackage[italian]{babel}
\usepackage[utf8]{inputenc}
\usepackage{amsmath}
\usepackage{booktabs}
\usepackage{graphicx}

% Stat block body: ridotto per far stare l'intera scheda in una pagina
\DndSetFonts[stat-block-body-style=\footnotesize]

\title{Obasi N'Dala \\
  \large Scheda del Personaggio --- Mago dell'Ordine degli Scriba}
\author{Giocatore}
\date{\today}

\begin{document}

\frontmatter
\maketitle
\tableofcontents
\mainmatter

% ------------------------------------------------------------
\part{Scheda}
% ------------------------------------------------------------

\chapter{Identità}

\chapter{Scheda del personaggio}

\begin{figure}[t]
  \centering
  \includegraphics[width=0.62\columnwidth]{ObasiNdala.jpeg}
  \caption{Obasi N'Dala, ritratto.}
\end{figure}

\begin{DndMonster}{Obasi N'Dala}
  \DndMonsterType{Umano Variante, Mago (Ordine degli Scriba) 5º, Caotico Buono}

  \DndMonsterBasics[
    armor-class = {16 (Armatura Magica + Mantello della Protezione)},
    hit-points  = {30 (8 + 4d6 + 4d6)},
    speed       = {9 m (30 ft)},
  ]

  \DndMonsterAbilityScores[
    str = 11,
    dex = 14,
    con = 14,
    int = 20,
    wis = 12,
    cha = 12,
  ]

  \DndMonsterDetails[
    saving-throws = {Int $+9$, Sag $+5$, Des $+3$, Cos $+3$, Car $+2$, For $+1$ (incl. +1 Mantello)},
    skills = {Arcano $+8$, Storia $+8$, Indagare $+8$, Religione $+8$, Natura $+5$},
    senses = {Percezione passiva 11},
    languages = {Corrente + una lingua a scelta},
    challenge = 5,
    proficiency-bonus, % = +3
  ]

  \DndMonsterSection{Tratti}

  \DndMonsterAction{Umano Variante (razza)}
  \textit{Incremento punteggi.} Due diverse caratteristiche a scelta
  ricevono $+1$. \textit{Abilità.} Competente in un'abilità a scelta.
  \textit{Talento.} Un talento a tua scelta (Resiliente COS).
  \textit{Taglia.} Media; velocità 9 m. \textit{Lingue.} Corrente e una
  lingua extra a scelta.

  \DndMonsterAction{Resiliente (COS)}
  $+1$ Cos; competente nei tiri salvezza su Costituzione.

  \DndMonsterAction{Recupero Arcano}
  Una volta al giorno, al termine di un riposo breve, recuperi slot
  incantesimo spesi per un livello complessivo $\le$ metà dei livelli da
  mago (arrotondato per eccesso), nessuno di 6º o superiore.

  \DndMonsterAction{Lancio Rituale}
  Puoi lanciare come rituale un incantesimo da mago con tag \emph{rituale}
  che sia nel tuo libro, senza averlo preparato.

  \DndMonsterAction{Libro Risvegliato}
  Usi il libro come focus incantatorio; puoi cambiare il tipo di danno di
  un incantesimo con uno di un altro incantesimo dello stesso livello nel
  libro; il lancio rituale usa il tempo normale invece di aggiungere 10
  minuti (1 volta al giorno).

  \DndMonsterAction{Penna Magica}
  Come azione bonus crei una penna che non necessita inchiostro; copiare
  un incantesimo richiede 2 minuti per livello; cancelli ciò che hai
  scritto sventolando la piuma (entro 1,5 m).

  \DndMonsterAction{Libro degli Incantesimi}
  Copiare un incantesimo nel libro richiede 2 ore e 50~mo per livello;
  copiare da una pergamena richiede una prova di Intelligenza (Arcano)
  con CD $10$ + livello (la pergamena viene distrutta in ogni caso).
  Preparare gli incantesimi del giorno richiede almeno 1 minuto per
  livello e si completa dopo un riposo lungo; puoi preparare un numero di
  incantesimi pari al tuo modificatore di Intelligenza + il tuo livello
  da mago (5 + 5 = 10).

  \DndMonsterAction{Ricercatore (Accolito)}
  Quando tenti di apprendere o ricordare un'informazione, se non la
  conosci sai spesso dove e da chi ottenerla: una biblioteca, uno
  scriptorium, un'università, un saggio o un'altra persona colta. Con il
  permesso del Dungeon Master puoi conoscere dettagli sulle religioni,
  le cerimonie e le divinità con cui hai a che fare.

  \DndMonsterSection{Azioni}

  \DndMonsterAction{Pugnale $+1$ (Mithril)}
  \textit{Mischia.} $+6$ a colpire, portata 1,5 m, un bersaglio. Colpito:
  4 (1d4$+3$) danni perforanti.

  \DndMonsterAction{Dardo di Fuoco (trucchetto)}
  \textit{A distanza incantata.} $+8$ a colpire, gittata 36 m, un
  bersaglio. Colpito: 2d10 danni da fuoco.

  \DndMonsterAction{Stretta Folgorante (trucchetto)}
  \textit{Mischia incantata.} $+8$ a colpire; vantaggio su bersagli con
  armatura di metallo. Colpito: 2d8 danni da fulmine e il bersaglio non
  può compiere reazioni fino all'inizio del tuo prossimo turno.

  \DndMonsterAction{Tocco Gelido (trucchetto)}
  \textit{A distanza incantata.} $+8$ a colpire, gittata 36 m, un
  bersaglio. Colpito: 1d8 danni necrotici ed entro la fine del prossimo
  turno del bersaglio, il primo tiro per danni di un attacco che gli
  infligge è svantaggiato.

  \DndMonsterAction{Dardo Incantato (slot 1º, 3) }
  \textit{Incantesimo.} Crei tre dardi (uno in più per ogni livello di
  slot oltre il 1º), ognuno infligge 1d4$+1$ danni da forza.

  \DndMonsterAction{Globo Cromatico (slot 1º) }
  \textit{A distanza incantata.} $+8$ a colpire, gittata 27 m. Colpito:
  3d8 (acido, freddo, fulmine, fuoco, tuono o veleno), +1d8 per slot oltre
  il 1º.

  \DndMonsterSection{Incantesimi (CD 16, bonus attacco $+8$)}

  \DndMonsterAction{Trucchetti (4)}
  Luce, Stretta Folgorante, Dardo di Fuoco, Tocco Gelido.

  \DndMonsterAction{Livello 1º (slot 4, 4 preparati)}
  Globo Cromatico, Armatura Magica, Dardo Incantato, Scudo.\\
  \emph{Rituali nel libro (non preparati):} Allarme, Identificare,
  Individuazione del Magico, Trova Famiglio.

  \DndMonsterAction{Livello 2º (slot 3, 4 preparati)}
  Allucinazione di Forza, Invisibilità, Soffio del Drago, Frusta Mentale
  di Tasha.

  \DndMonsterAction{Livello 3º (slot 2, 2 preparati)}
  Velocità, Palla di Fuoco (8d6, TS Des CD 16, sfera 6 m).

  \DndMonsterSection{Equipaggiamento}
  Pugnale $+1$ (mithril), Borsa dei Componenti, Zaino, Penna, 10 fogli di
  pergamena, sacchetto di sabbia, coltello, abito comune, 12 pergamene per
  incantesimi.

  \DndMonsterSection{Oggetti magici}
  Anello Accumula Incantesimi (raro), Perla del Potere (non comune),
  Mantello della Protezione (non comune, $+1$ CA e TS), Buco Portatile
  (raro).

  Carico max 82,5~kg -- Monete: 708~mo.
\end{DndMonster}

% ------------------------------------------------------------
\part{Equipaggiamento e Magie}
% ------------------------------------------------------------

\chapter{Equipaggiamento}
\section{Oggetti ordinari}
\begin{itemize}
  \item 1 Pugnale $+1$ con elsa in bronzo e lama in mithril
  \item 1 Borsa dei Componenti
  \item 1 Zaino
  \item 1 Penna
  \item 10 Fogli di pergamena (non per incantesimi)
  \item 1 Sacchetto di sabbia
  \item 1 Coltello
  \item 1 Abito comune
  \item 12 Pergamene per incantesimi
\end{itemize}

Capacità di carico: \textbf{82,5~kg}. Monete: \textbf{708~mo}.

\section{Oggetti magici}
\begin{DndTable}[header=Inventario magico]{cl}
  Anello Accumula Incantesimi & Raro \\
  Perla del Potere            & Non Comune \\
  Mantello della Protezione   & Non Comune \\
  Buco Portatile              & Raro \\
\end{DndTable}

\subsection*{Anello Accumula Incantesimi} (Raro, sintonizzazione)
Questo anello immagazzina gli incantesimi lanciati al suo interno,
conservandoli finché chi vi è in sintonia non li usa. L'anello può
conservare fino a 5 livelli di incantesimi alla volta. Qualsiasi creatura
può lanciare un incantesimo dal 1\textsuperscript{o} al 5\textsuperscript{o}
livello sull'anello, toccandolo mentre lo lancia: l'incantesimo non ha
effetto, se non essere immagazzinato. Il livello dello slot usato determina
quanto spazio occupa. Mentre indossi l'anello, puoi lanciare qualunque
incantesimo contenga: usa livelli slot, CD del tiro salvezza, bonus
d'attacco e caratteristica del lanciatore originario, ma per il resto è
considerato come se lo lanciassi tu. L'incantesimo lanciato non è più
contenuto nell'anello, liberando spazio.

\subsection*{Perla del Potere} (Non Comune, sintonizzazione)
Mentre hai la perla sulla tua persona, puoi usare un'azione per pronunciare
la sua parola di comando e recuperare uno slot incantesimo speso. Se lo slot
speso era di 4\textsuperscript{o} livello o superiore, il nuovo slot è di
3\textsuperscript{o} livello. Una volta usata, la perla non può essere
riusata fino all'alba successiva.

\subsection*{Mantello della Protezione} (Non Comune, sintonizzazione)
Mentre indossi questo mantello, ottieni un bonus di $+1$ alla Classe Armatura
e ai tiri salvezza.

\subsection*{Buco Portatile} (Raro)
Questo sottile panno nero, morbido come la seta, è piegato fino alle
dimensioni di un fazzoletto. Si apre in un disco circolare di 1,8 metri di
diametro. Con un'azione lo apri e lo posizioni su una superficie solida, dove
crea un buco extradimensionale profondo 3 metri. Lo spazio cilindrico esiste
su un piano diverso, quindi non può creare passaggi aperti. Anche se pieno,
il buco non pesa praticamente nulla. Con un'azione lo richiudi ripiegandolo;
creature o oggetti al suo interno rimangono nello spazio
extradimensionale. Se richiuso con qualcuno dentro, questa può con un'azione
compiere una prova di Forza CD 10 per farsi largo e uscire entro 1,5 metri
dal buco o da chi lo trasporta. Una creatura che respira dentro un buco
chiuso può sopravvivere fino a 10 minuti, poi inizia a soffocare. Inserire
il buco portatile in uno spazio extradimensionale (come quello di una borsa
conservante) distrugge istantaneamente entrambi gli oggetti e apre un
portale verso il Piano Astrale.

\section{Attacchi}
\begin{DndTable}[header=Attacchi]{cccc}
  \textbf{Attacco} & \textbf{Bonus} & \textbf{Danno/Tipo} & \textbf{Note} \\
  Pugnale $+1$ (mithril) & $+6$ & 1d4$+3$ Perforante & -- \\
  Dardo di Fuoco & $+8$ & 2d10 Fuoco & Trucchetto \\
  Stretta Folgorante & $+8$ & 2d8 Fulmine & Trucchetto \\
  Tocco Gelido & $+8$ & 1d8 Necrotici & Trucchetto, svantaggio al tiro per danni \\
  Globo Cromatico & $+8$ & 3d8 Elementale & Slot 1 (o superiore) \\
  Dardo Incantato & $+8$ & 3d4$+3$ Forza & Slot 1 (o superiore), multi bersaglio \\
  Soffio del Drago & -- & 3d6 Elementale & TSD Destrezza, cono \\
  Frusta Mentale di Tasha & -- & 3d6 Psichici & TS Intelligenza, CD 16 \\
  Allucinazione di Forza & -- & 1d6 Psichici/rnd & TS Intelligenza \\
  Palla di Fuoco & -- & 8d6 Fuoco & TS Destrezza, sfera \\
\end{DndTable}

\chapter{Incantesimi}
\section{Trucchetti (4)}
Luce \textbullet{} Stretta Folgorante \textbullet{} Dardo di Fuoco
\textbullet{} Tocco Gelido

\section{Livello 1 (slot: 4) -- 4 preparati}
Globo Cromatico \textbullet{} Armatura Magica \textbullet{} Dardo Incantato
\textbullet{} Scudo\\
\emph{Rituali nel libro (non preparati):} Allarme \textbullet{} Identificare
\textbullet{} Individuazione del Magico \textbullet{} Trova Famiglio

\section{Livello 2 (slot: 3) -- 4 preparati}
Allucinazione di Forza \textbullet{} Invisibilità \textbullet{}
Soffio del Drago \textbullet{} Frusta Mentale di Tasha

\section{Livello 3 (slot: 2) -- 2 preparati}
Velocità \textbullet{} Palla di Fuoco

% ------------------------------------------------------------
\part{Background}
% ------------------------------------------------------------

\chapter{Il personaggio}

\DndDropCapLine{O}{basi N'Dala} è un mago dell'Ordine degli Scriba,
cresciuto originariamente come \emph{Accolito}; la sua mente è
un'arma affilata quanto la sua penna.

\section{Biografia}
Nato nei bassifondi, Obasi ha sempre saputo che la magia non è un dono
mistico, ma una scienza esatta, un'arma che richiede disciplina e
spietatezza. È, in segreto, \textbf{un maestro del culto di Mystra}, la
Madre di Tutta la Magia: per lui la Trama non è solo uno strumento, ma il
tessuto sacro dell'universo. Il suo percorso iniziò quando incrociò la
strada del \emph{Maestro Scriba}, un anziano e cinico mago che operava
nell'ombra, segretamente affiliato alla Rete Nera (Zhentarim). Fu lui a
insegnargli che la Trama è un arazzo perfetto, e che chiunque vi acceda
senza studio --- come quegli abomini degli stregoni, che vomitano magia
per puro caso genetico --- è una pericolosa anomalia che va studiata,
ingabbiata e neutralizzata. Solo i maghi meritano il potere.

Sotto la guida del Maestro, Obasi divenne un membro dell'Ordine degli
Scriba. Imparò a incanalare la sua essenza magica in un oggetto
apparentemente innocuo. Obasi, con il suo aspetto rilassato, i lunghi
rasta brizzolati e la voce sempre bassa e pacata, scelse un semplice
fermacapelli d'osso scuro che teneva raccolti i suoi dreadlock. Basta un
gesto fluido della sua mano perché quell'ornamento si sfili dai capelli
e, brillando di un'energia spettrale, si trasformi nella sua Penna del
Mago, strumento di creazione e di morte.

La sua flemma e la sua spietatezza attirarono l'attenzione dei vertici
della Rete Nera, che lo infiltrarono nei ranghi dei mercenari magici di
Waterdeep. Obasi recitò la sua parte alla perfezione: un mago taciturno,
utile ma apparentemente privo di ambizioni politiche. Un profilo tanto
perfetto da attirare l'attenzione di Vajra Safahr, l'attuale Bastone
Nero.

Vajra, tuttavia, non è sprovveduta. Vide in Obasi un individuo fin
troppo capace e potenzialmente pericoloso, il classico ``soggetto
indesiderabile''. Non sospettando affatto la sua affiliazione agli
Zhentarim, decise di usarlo come carne da macello: lo convocò in segreto
e gli affidò una missione suicida, infiltrarsi nel covo di Xanathar per
sgominare il Beholder che controllava i vicoli di Waterdeep. Obasi
accettò con un cenno del capo --- non per obbedire al Bastone Nero, ma
perché sapeva che tra i trofei del Beholder giaceva il \textbf{filatterio
di un lich}, nascosto in un anello posseduto da Xanathar, che la Rete
Nera e i suoi maestri del culto di Mystra volevano recuperare a ogni
costo.

Contro ogni aspettativa del Bastone Nero, Obasi ebbe successo: la gilda
fu disgregata e il filatterio ritornò nelle mani giuste. Si era alla
primavera del 1492 DR, poco dopo il \emph{Furto dei Draghi}, e già i
vicoli di Waterdeep tremavano per la guerra tra la gilda di Xanathar e la
Rete Nera. Quando Obasi tornò, con la solita flemma e la voce priva di
affanno, Vajra capì di avere un problema ancora più grosso tra le mani.
Con un finto sorriso di compiacimento, il Bastone Nero gli offrì una
``promozione'': la missione definitiva per un eroe, scendere a Sottomonte,
trovare il folle Halaster Blackcloak e liberare Waterdeep dalla sua
minaccia per sempre.

Per questa impresa Obasi radunò una squadra: \textbf{Yelena}, una rogue
che con disinvoltura mescola pugnali e incantesimi --- l'ha vista
lanciare magia e ne è rimasto affascinato, e tra loro è nato un buon
rapporto; lei, inoltre, lo sta addestrando all'uso della lama;
\textbf{Varis}, un arciere eccezionale che gli ha salvato la vita più di
una volta e gode della sua piena fiducia; \textbf{Opa}, un chierico della
vita nano che lo sta istruendo all'uso dello scudo; e \textbf{Theren}, un
umano di cui sa poco e di cui --- per istinto --- non si fida del tutto.

Vajra è convinta di averlo mandato a morire nel più grande tritacarne di
Faerûn. Non sa che Obasi ha sorriso a quella proposta. La missione si
sposa perfettamente con il suo piano e con gli interessi della Rete Nera:
scenderà a Sottomonte, sì. Ucciderà Halaster, certo. Ma non per liberare
Waterdeep. Perché Sottomonte, intrisa com'è di magia fino al midollo, sta
prosciugando la Trama in tutta questa parte di Faerûn --- e se il drenaggio
continua, la Trama stessa si lacererà e \textbf{Mystra morirà}. Obasi
scenderà per uccidere il Mago Pazzo e ripristinare il tessuto della Trama,
poi reclamare Sottomonte per sé, impossessandosi del gigantesco nodo
magico che vi si nasconde, per risalire in superficie, schiacciare il
Bastone Nero, rovesciare le istituzioni che limitano i maghi e instaurare
un nuovo ordine dove la magia è controllata solo da chi ha il diritto di
farlo.

\section{Tratti caratteriali}
Parlo sempre con un tono di voce basso, calmo e rilassato, anche mentre
scateno magie letali. La mia espressione pacifica e il mio aspetto
trasandato nascondono una mente calcolatrice e spietata.

\section{Ideali}
La magia è la forza suprema dell'universo e deve essere dominata solo da
chi ha la disciplina per studiarla. Chi nasce con poteri innati è
un'anomalia della natura e deve essere messo in gabbia o eliminato.

\section{Legami}
Il mio pennino-fermacapelli è il mio compagno più fidato, l'estensione
della mia mente. La Rete Nera è lo strumento con cui porterò l'ordine a
Waterdeep e mi prenderò Sottomonte.

\section{Difetti}
La mia superbia intellettuale è immensa: sottovaluto sistematicamente
stregoni e incantatori innati, considerandoli animali indisciplinati. In
più, la mia ambizione di dominare Sottomonte mi rende cieco di fronte
alla vera, insondabile follia di quel luogo.

\begin{DndComment}{Note del giocatore}
Riserva questo riquadro per annotazioni, aggiornamenti di livello,
equipaggiamento o dettagli della campagna.
\end{DndComment}

\end{document}
